\chapter{Chapter06 — Finite Fields}
%%% generated by the blueprint pipeline from FABL.Chapter06.Constructions.FiniteFields
\section{Overview}

Declarations of the Lean module \texttt{FABL.Chapter06.Constructions.FiniteFields}.

\section{Declarations}

\begin{definition}[Binary extension field]
\lean{FABL.BinaryExtensionField}
\statementsource{Vad12}{
  pdf-bf2ce5422088-p014-b002,
  pdf-bf2ce5422088-p068-b007,
  pdf-bf2ce5422088-p148-b001
}
\statementsource{Analysis-of-Boolean-Functions-by-Ryan-ODonnell}{
  pdf-ffbe9c7b3306-p016-b002,
  pdf-ffbe9c7b3306-p153-b001
}
\statementsource{ghsz-combinatorial-list-decoding}{pdf-6bd33ae73ade-p002-b001}
\label{FABL.BinaryExtensionField}
\leanok
For a natural number $\ell$, the binary extension field of degree $\ell$ is the finite
field $\mathbb{F}_{2^\ell}$.
\end{definition}

\begin{theorem}[Dimension of the binary extension field]
\lean{FABL.binaryExtensionField_finrank}
\label{FABL.binaryExtensionField_finrank}
\leanok
\uses{FABL.BinaryExtensionField, FABL.𝔽₂}
\difficulty{1}
For every positive integer $\ell$, the degree-$\ell$ extension field of $\mathbb{F}_2$
has dimension $\ell$ as a vector space over $\mathbb{F}_2$.
\end{theorem}

\begin{theorem}[Binary extension field nat card]
\lean{FABL.binaryExtensionField_natCard}
\statementsource{Vad12}{
  pdf-bf2ce5422088-p068-b007,
  pdf-bf2ce5422088-p148-b001
}
\statementsource{Analysis-of-Boolean-Functions-by-Ryan-ODonnell}{
  pdf-ffbe9c7b3306-p016-b002,
  pdf-ffbe9c7b3306-p153-b001
}
\label{FABL.binaryExtensionField_natCard}
\leanok
\uses{FABL.BinaryExtensionField}
\difficulty{1}
% dataset-pipeline: stub -- prose pending build_informal_statements.py
Statement pending regeneration.
\end{theorem}

\begin{definition}[Basis of the binary extension field]
\lean{FABL.binaryExtensionBasis}
\statementsource{Analysis-of-Boolean-Functions-by-Ryan-ODonnell}{pdf-ffbe9c7b3306-p153-b001}
\label{FABL.binaryExtensionBasis}
\leanok
\uses{FABL.BinaryExtensionField, FABL.binaryExtensionField_finrank, FABL.𝔽₂}
Let $\ell$ be a nonzero natural number. The binary extension field of degree $\ell$ has
a basis over $\mathbb{F}_2$ indexed by $\{0,\ldots,\ell-1\}$.
\end{definition}

\begin{definition}[Binary extension encoding]
\lean{FABL.binaryExtensionEncode}
\statementsource{Analysis-of-Boolean-Functions-by-Ryan-ODonnell}{pdf-ffbe9c7b3306-p153-b001}
\label{FABL.binaryExtensionEncode}
\leanok
\uses{FABL.BinaryExtensionField, FABL.F₂Cube, FABL.binaryExtensionBasis, FABL.𝔽₂}
Let $\ell$ be a nonzero natural number. There is a linear equivalence over
$\mathbb{F}_2$ between the binary extension field of degree $\ell$ and the space of
length-$\ell$ vectors over $\mathbb{F}_2$.
\end{definition}

\begin{theorem}[Encoding zero]
\lean{FABL.binaryExtensionEncode_zero}
\proofsource{Analysis-of-Boolean-Functions-by-Ryan-ODonnell}{pdf-ffbe9c7b3306-p153-b001}
\label{FABL.binaryExtensionEncode_zero}
\leanok
\uses{FABL.binaryExtensionEncode}
\proofstep
  {FABL.𝔽₂}
  {direct}
  {Analysis-of-Boolean-Functions-by-Ryan-ODonnell}
  {pdf-ffbe9c7b3306-p015-b008}
\proofstep
  {FABL.F₂Cube}
  {context}
  {Analysis-of-Boolean-Functions-by-Ryan-ODonnell}
  {pdf-ffbe9c7b3306-p015-b008,
    pdf-ffbe9c7b3306-p031-b002}
\proofstep
  {FABL.Sign}
  {context}
  {Analysis-of-Boolean-Functions-by-Ryan-ODonnell}
  {pdf-ffbe9c7b3306-p014-b002,
    pdf-ffbe9c7b3306-p022-b004}
\proofstep
  {FABL.signValue}
  {context}
  {Analysis-of-Boolean-Functions-by-Ryan-ODonnell}
  {pdf-ffbe9c7b3306-p014-b002,
    pdf-ffbe9c7b3306-p022-b004}
\proofstep
  {FABL.signValue_one}
  {context}
  {Analysis-of-Boolean-Functions-by-Ryan-ODonnell}
  {pdf-ffbe9c7b3306-p014-b002}
\proofstep
  {FABL.signValue_neg_one}
  {context}
  {Analysis-of-Boolean-Functions-by-Ryan-ODonnell}
  {pdf-ffbe9c7b3306-p014-b002}
\proofstep
  {FABL.monomial}
  {direct}
  {Analysis-of-Boolean-Functions-by-Ryan-ODonnell}
  {pdf-ffbe9c7b3306-p022-b001,
    pdf-ffbe9c7b3306-p022-b004}
\proofstep
  {FABL.fourierCoeff}
  {direct}
  {Analysis-of-Boolean-Functions-by-Ryan-ODonnell}
  {pdf-ffbe9c7b3306-p025-b004}
\proofstep
  {FABL.fourierSupport}
  {context}
  {Analysis-of-Boolean-Functions-by-Ryan-ODonnell}
  {pdf-ffbe9c7b3306-p022-b002,
    pdf-ffbe9c7b3306-p018-b003}
\proofstep
  {FABL.mem_fourierSupport}
  {context}
  {Analysis-of-Boolean-Functions-by-Ryan-ODonnell}
  {pdf-ffbe9c7b3306-p022-b002}
\proofstep
  {FABL.granularityFrequencyWithin}
  {context}
  {Analysis-of-Boolean-Functions-by-Ryan-ODonnell}
  {pdf-ffbe9c7b3306-p076-b001,
    pdf-ffbe9c7b3306-p076-b003}
\proofstep
  {FABL.mem_granularityFrequencyWithin}
  {context}
  {Analysis-of-Boolean-Functions-by-Ryan-ODonnell}
  {pdf-ffbe9c7b3306-p076-b001,
    pdf-ffbe9c7b3306-p076-b003}
\proofstep
  {FABL.BinaryExtensionField}
  {context}
  {Analysis-of-Boolean-Functions-by-Ryan-ODonnell}
  {pdf-ffbe9c7b3306-p016-b002,
    pdf-ffbe9c7b3306-p153-b001}
\proofstep
  {FABL.binaryExtensionField_finrank}
  {context}
  {Analysis-of-Boolean-Functions-by-Ryan-ODonnell}
  {pdf-ffbe9c7b3306-p153-b001}
\proofstep
  {FABL.binaryExtensionBasis}
  {context}
  {Analysis-of-Boolean-Functions-by-Ryan-ODonnell}
  {pdf-ffbe9c7b3306-p153-b001}
\proofstep
  {FABL.binaryExtensionEncode}
  {context}
  {Analysis-of-Boolean-Functions-by-Ryan-ODonnell}
  {pdf-ffbe9c7b3306-p153-b001}
\proofstep
  {FABL.binaryExtensionEncode_add}
  {context}
  {Analysis-of-Boolean-Functions-by-Ryan-ODonnell}
  {pdf-ffbe9c7b3306-p153-b001}
\difficulty{1}
Let $\ell$ be a positive integer. The linear encoding of the degree-$\ell$ binary
extension field sends its zero element to the zero vector.
\end{theorem}

\begin{theorem}[Binary extension encode add]
\lean{FABL.binaryExtensionEncode_add}
\statementsource{Analysis-of-Boolean-Functions-by-Ryan-ODonnell}{pdf-ffbe9c7b3306-p153-b001}
\label{FABL.binaryExtensionEncode_add}
\leanok
\uses{FABL.BinaryExtensionField, FABL.binaryExtensionEncode}
\difficulty{1}
% dataset-pipeline: stub -- prose pending build_informal_statements.py
Statement pending regeneration.
\end{theorem}

\begin{theorem}[Root count and degree]
\lean{FABL.ncard_rootSet_le_natDegree}
\proofsource{Vad12}{
  pdf-bf2ce5422088-p017-b002,
  pdf-bf2ce5422088-p044-b001
}
\proofstep
  {FABL.𝔽₂}
  {context}
  {Vad12}
  {pdf-bf2ce5422088-p164-b005}
\proofstep
  {FABL.F₂Cube}
  {context}
  {Vad12}
  {pdf-bf2ce5422088-p142-b003}
\proofstep
  {FABL.Sign}
  {context}
  {Vad12}
  {pdf-bf2ce5422088-p308-b005}
\proofstep
  {FABL.signValue}
  {context}
  {Vad12}
  {pdf-bf2ce5422088-p308-b005}
\proofstep
  {FABL.signValue_one}
  {context}
  {Vad12}
  {pdf-bf2ce5422088-p308-b005}
\proofstep
  {FABL.signValue_neg_one}
  {context}
  {Vad12}
  {pdf-bf2ce5422088-p308-b005}
\proofstep
  {FABL.monomial}
  {context}
  {Vad12}
  {pdf-bf2ce5422088-p308-b005}
\proofstep
  {FABL.fourierCoeff}
  {context}
  {Vad12}
  {pdf-bf2ce5422088-p308-b005}
\proofstep
  {FABL.fourierSupport}
  {context}
  {Vad12}
  {pdf-bf2ce5422088-p308-b005}
\proofstep
  {FABL.mem_fourierSupport}
  {context}
  {Vad12}
  {pdf-bf2ce5422088-p308-b005}
\proofstep
  {FABL.granularityFrequencyWithin}
  {context}
  {Vad12}
  {pdf-bf2ce5422088-p308-b005}
\proofstep
  {FABL.mem_granularityFrequencyWithin}
  {context}
  {Vad12}
  {pdf-bf2ce5422088-p308-b005}
\proofstep
  {FABL.BinaryExtensionField}
  {context}
  {Vad12}
  {pdf-bf2ce5422088-p069-b001}
\proofstep
  {FABL.binaryExtensionField_finrank}
  {context}
  {Vad12}
  {pdf-bf2ce5422088-p069-b001}
\proofstep
  {FABL.binaryExtensionBasis}
  {context}
  {Vad12}
  {pdf-bf2ce5422088-p069-b001}
\proofstep
  {FABL.binaryExtensionEncode}
  {context}
  {Vad12}
  {pdf-bf2ce5422088-p069-b001}
\proofstep
  {FABL.binaryExtensionEncode_zero}
  {context}
  {Vad12}
  {pdf-bf2ce5422088-p069-b001}
\proofstep
  {FABL.binaryExtensionEncode_add}
  {context}
  {Vad12}
  {pdf-bf2ce5422088-p069-b001}
\label{FABL.ncard_rootSet_le_natDegree}
\leanok
\difficulty{1}
Let $K$ be a field and let $p\in K[X]$. The number of distinct roots of $p$ in $K$ is at
most the natural degree of $p$.
\end{theorem}

\begin{theorem}[Nonvanishing of the Vandermonde determinant]
\lean{FABL.det_vandermonde_ne_zero_of_injective}
\label{FABL.det_vandermonde_ne_zero_of_injective}
\leanok
\difficulty{1}
Let $K$ be a field, let $k$ be a nonnegative integer, and let $\alpha_1,\ldots,\alpha_k$
be pairwise distinct elements of $K$. The determinant of the Vandermonde matrix formed
from these elements is nonzero.
\end{theorem}

\begin{theorem}[Kernel of a Vandermonde matrix]
\lean{FABL.eq_zero_of_vandermonde_mulVec_eq_zero}
\label{FABL.eq_zero_of_vandermonde_mulVec_eq_zero}
\leanok
\uses{FABL.det_vandermonde_ne_zero_of_injective}
\difficulty{1}
Let $K$ be a field, let $k$ be a natural number, and let $\alpha_1,\ldots,\alpha_k$ be
pairwise distinct elements of $K$. If a vector $v\in K^k$ is in the kernel of the
Vandermonde matrix formed from these elements, then $v=0$.
\end{theorem}
